\rtmtight
\begin{tblr}{width=\textwidth,colspec={|Q[c,m,2.1cm]|X[l,m]|},hlines,vlines,hspan=minimal,rowsep=1.5pt,colsep=3pt,row{1}={bg=headgray,font=\bfseries}}
\SetCell{c,m}\hfil\logo\hfil & \textbf{POLITEKNIK NEGERI MALANG}\par\textbf{JURUSAN TEKNOLOGI INFORMASI}\par\textbf{PROGRAM STUDI: D4 SISTEM INFORMASI BISNIS} \\
\SetCell[c=2]{l} \textbf{RENCANA TUGAS MAHASISWA (RTM) DAN RUBRIK PENILAIAN} & \\
Nama Mata Kuliah & Pemrograman Web Lanjut \\
Kode Mata Kuliah & RTI254007 \quad \textbf{sks / Jam perminggu:} 3 SKS/6 jam \quad \textbf{Semester:} 4 \\
Dosen Pengampu & \begin{enumerate}\item Dian Hanifudin Subhi, S.Kom., M.Kom.\item Ade Ismail, S.Kom., M.TI.\end{enumerate} \\
\end{tblr}

\assessmentblock{Tugas Praktik: Framework, Routing, \& Basis Data}{Hands-on Practice}{SCPMK0202-04001, SCPMK0202-04002}{Mahasiswa menyelesaikan tugas praktik mingguan terkait setup framework, routing, controller, desain basis data, migrasi skema, dan server-side templating; setiap tugas merupakan increment yang membangun aplikasi studi kasus POS/kasir UMKM secara bertahap.}{Menganalisis kebutuhan, mengerjakan artefak praktik, menguji hasil, mendokumentasikan evidence, dan mempresentasikan keputusan bila diminta.}{Laporan, artefak digital, repository/prototype, evidence pengujian, dan/atau bahan presentasi.}{Kesiapan lingkungan \& konfigurasi framework & 14 & Framework terpasang, konfigurasi basis data, struktur proyek, dan version control siap digunakan secara penuh. & Framework berjalan tetapi konfigurasi atau struktur awal belum lengkap. & Framework tidak berjalan atau konfigurasi dasar salah. \\ \\
Routing, controller, basis data, \& migrasi skema & 12.3 & Route, controller, migrasi skema, model, dan relasi dasar sesuai kebutuhan modul serta konsisten. & Komponen utama tersedia tetapi relasi, migrasi, atau controller belum sepenuhnya konsisten. & Routing/model/migrasi tidak terhubung atau tidak sesuai fitur. \\ \\
Templating, validasi, \& evidence praktik & 8.7 & Server-side template, formulir, validasi, feedback error, dan screenshot/commit praktik terdokumentasi jelas. & Template dan formulir berjalan tetapi validasi atau evidence praktik masih terbatas. & Template tidak mendukung alur fitur atau tidak ada bukti praktik. \\}

\assessmentblock{UTS: Evaluasi Milestone Aplikasi Studi Kasus}{UTS praktik}{SCPMK0202-04001, SCPMK0202-04002, SCPMK0202-04003}{Mahasiswa mendemonstrasikan milestone aplikasi studi kasus POS/kasir UMKM (hasil increment minggu 1--7) yang mencakup CRUD lengkap, autentikasi, otorisasi, dan arsitektur MVC, kemudian mempertahankan keputusan teknis secara lisan (viva). Evaluasi bersifat individu, bukan proyek kelompok PBL.}{Menganalisis kebutuhan, mengerjakan artefak proyek, menguji hasil, mendokumentasikan evidence, dan mempertahankan keputusan teknis secara lisan.}{Aplikasi web berjalan, repository, evidence pengujian, dan/atau bahan viva.}{Kelengkapan CRUD \& validasi & 8 & Create, read, update, delete, validasi input, flash message, dan penanganan error berjalan lengkap. & CRUD utama berjalan tetapi validasi atau penanganan error belum lengkap. & CRUD tidak berjalan utuh atau data tidak tersimpan benar. \\ \\
Kualitas arsitektur MVC \& basis data & 7 & Route, controller, model, migrasi, seeder, dan view dipisahkan rapi mengikuti konvensi framework yang digunakan. & Struktur cukup mengikuti arsitektur MVC tetapi masih ada logika bercampur atau konvensi tidak konsisten. & Kode tidak mengikuti arsitektur MVC atau sulit dipelihara. \\ \\
Pengujian praktik \& defense & 5 & Skenario CRUD normal/invalid diuji, data diverifikasi di basis data, dan keputusan teknis dapat dijelaskan secara lisan. & Pengujian dilakukan terbatas atau defense belum menjelaskan semua keputusan. & Tidak ada pengujian bermakna atau tidak mampu menjelaskan modul. \\}

\assessmentblock{Tugas Praktik: Autentikasi, API, \& Pengolahan Data}{Hands-on Practice}{SCPMK0202-04003, SCPMK0202-04004}{Mahasiswa menyelesaikan tugas praktik terkait autentikasi, otorisasi berbasis peran, desain API, token-based authentication, konsumsi API dari sisi klien, serta pengolahan dan ekspor data; setiap tugas merupakan increment lanjutan yang menyelesaikan aplikasi studi kasus POS/kasir UMKM.}{Menganalisis kebutuhan, mengerjakan artefak praktik, menguji hasil, mendokumentasikan evidence, dan mempresentasikan keputusan bila diminta.}{Laporan, artefak digital, repository/prototype, evidence pengujian, dan/atau bahan presentasi.}{Autentikasi \& otorisasi berbasis peran & 10 & Login, logout, manajemen peran/hak akses, proteksi route, dan akses data diterapkan sesuai skenario. & Autentikasi berjalan tetapi manajemen peran, proteksi route, atau akses data masih memiliki celah. & Autentikasi/otorisasi tidak berjalan atau route penting tidak terlindungi. \\ \\
Desain API \& pengolahan data & 5.8 & Endpoint API, request/response, validasi, penanganan error, import/export data, dan token-based auth bekerja konsisten. & Fitur API/pengolahan data berjalan sebagian tetapi validasi atau penanganan error belum lengkap. & API gagal atau fitur import/export tidak sesuai kontrak data. \\ \\
Konsumsi API dari sisi klien & 3 & Contoh klien (consumer) yang mengonsumsi API bekerja dengan benar: berhasil login, menyimpan/mengirim token, memanggil endpoint terproteksi, dan menangani respons/galat dengan tepat. & Klien berhasil memanggil sebagian endpoint tetapi penanganan token atau galat belum konsisten. & Klien tidak berhasil mengonsumsi API atau tidak ada contoh klien sama sekali. \\ \\
Dokumentasi \& evidence pengujian & 6.2 & Dokumentasi endpoint (OpenAPI/Swagger atau setara), contoh request/response, dan evidence eksekusi tersedia jelas. & Dokumentasi tersedia tetapi cakupan dan evidence masih terbatas. & Tidak ada dokumentasi memadai atau evidence tidak membuktikan fitur. \\}

\assessmentblock{PBL: Proyek Aplikasi Web Terintegrasi}{Project Based Learning}{SCPMK0202-04004}{Mahasiswa mengembangkan proyek aplikasi web berbasis framework secara terintegrasi dalam tim pada domain bisnis pilihan tim (mis. inventori, koperasi, reservasi, laundry), terinspirasi dari aplikasi studi kasus POS/kasir UMKM, mencakup perencanaan, pengembangan fitur, integrasi modul, optimasi, dan finalisasi sesuai kebutuhan pengguna.}{Menganalisis kebutuhan pengguna, merencanakan arsitektur, mengerjakan artefak proyek secara bertahap, menguji hasil, mendokumentasikan evidence, dan mempresentasikan produk.}{Laporan proyek, artefak digital, repository, video demo, dan/atau bahan presentasi.}{Kesesuaian kebutuhan \& rancangan aplikasi & 4 & Kebutuhan pengguna, model data, peran, dan prioritas fitur dirancang konsisten dan terdokumentasi. & Rancangan cukup sesuai tetapi sebagian kebutuhan atau peran belum terpetakan jelas. & Aplikasi tidak menunjukkan kebutuhan pengguna atau rancangan tidak konsisten. \\ \\
Integrasi fitur \& kualitas implementasi & 3.5 & CRUD, autentikasi, API/import-export, validasi, relasi basis data, dan antarmuka terintegrasi secara stabil. & Fitur utama terintegrasi tetapi masih ada kekurangan validasi, relasi, atau penanganan error. & Fitur utama tidak terintegrasi atau aplikasi sering gagal. \\ \\
Pengujian, deployment, \& presentasi proyek & 2.5 & Pengujian, data demo/seeder, deployment atau demo lokal, repository, dan presentasi keputusan teknis tersedia lengkap. & Evidence proyek tersedia tetapi pengujian/deployment atau presentasi keputusan belum kuat. & Tidak ada evidence memadai atau proyek tidak dapat didemonstrasikan. \\}

\assessmentblock{UAS: Evaluasi Proyek Aplikasi Web Final}{UAS praktik dan defense}{SCPMK0202-04004}{Mahasiswa mendemonstrasikan proyek web final, mempertahankan keputusan teknis secara lisan, dan mengevaluasi ketercapaian kebutuhan pengguna.}{Menganalisis ketercapaian kebutuhan, mendemonstrasikan aplikasi, mempertahankan keputusan teknis, dan memberikan rekomendasi peningkatan.}{Aplikasi web final, repository, dokumentasi teknis, dan/atau bahan presentasi/defense.}{Evaluasi kelengkapan fitur akhir & 4 & Fitur aplikasi memenuhi scope, peran, alur kebutuhan pengguna, validasi, dan pengelolaan data secara menyeluruh. & Fitur utama memenuhi scope tetapi beberapa validasi atau peran masih belum lengkap. & Fitur akhir tidak memenuhi scope atau banyak alur gagal. \\ \\
Kualitas teknis \& maintainability & 3.5 & Kode mengikuti konvensi framework, struktur rapi, query efisien, dan konfigurasi siap dijalankan ulang. & Kode cukup rapi tetapi masih ada duplikasi, query kurang efisien, atau konfigurasi kurang jelas. & Kode sulit dipelihara atau tidak dapat dijalankan ulang. \\ \\
Defense, evidence, \& rekomendasi & 2.5 & Defense menjelaskan arsitektur, pengujian, deployment, risiko, dan rekomendasi peningkatan aplikasi secara jelas. & Defense cukup jelas tetapi evidence pengujian/deployment atau rekomendasi masih terbatas. & Tidak mampu mempertahankan keputusan teknis atau evidence tidak memadai. \\}

\begin{tblr}{width=\textwidth,colspec={|Q[l,m,3.2cm]|X[l,m]|},hlines,vlines,hspan=minimal,row{1-3}={bg=headgray},rowsep=1.5pt,colsep=3pt}
Jadwal Pelaksanaan: & Minggu 1--16 sesuai RPS. Setiap evaluasi memiliki RTM dan rubrik multi-kriteria. \\
Lain-Lain Yang Diperlukan: & LMS, repositori/prototype/proyek, perangkat lunak sesuai mata kuliah, dan perangkat presentasi. \\
Pustaka: & Mengacu pustaka utama dan pendukung pada bagian RPS. \\
\end{tblr}
